% Prozentrechnung · Prozentwert · Prüfungs-Fokus MSA – bank/prozentrechnung/e3.jsonl (zusammenbau v0.6, Prüfungs-Fokus)
\einheitenkopf[e3]{Einheit 3 · Prozentwert berechnen}
\zweigzeile{ab Kl. 7 · P10 ×7}
\setcounter{aufgabe}{0}

% Anlauf (ohne Prüfkennung)
\begin{aufgabe}{Prozentwert – Anlauf}
\begin{teile}
\teil $50\,\%$ von $80\,€$ – wie viel?

\streifen[2]{50}{$0$}{$80\,€$}
\leerfeld[€]
\teil $6\,\%$ von $400$ kg \leerfeld[kg]
\end{teile}
\end{aufgabe}

\begin{aufgabe}{Prozentwert – Prüfungsaufgaben}
\begin{teile}
\teil Ein Fernseher kostet $480\,€$. Bei Barzahlung gibt es $15\,\%$ Rabatt. Wie viel Euro spart man? Kreuze an. \hfill (P21)\\ \kreuz{$32\,€$}\\ \kreuz{$408\,€$}\\ \kreuz{$72\,€$}\\ \kreuz{$552\,€$}
\teil Eine Waschmaschine kostet $650\,€$. Im Angebot gibt es $30\,\%$ Rabatt. Wie viel Euro spart man? Kreuze an. \hfill (P21)\\ \kreuz{$455\,€$}\\ \kreuz{$65\,€$}\\ \kreuz{$845\,€$}\\ \kreuz{$195\,€$}
\teil Ein Tablet kostet $240{,}00\,€$. Es gibt $15\,\%$ Rabatt. Wie viel Euro Rabatt sind das? \hfill (P17) \\ \leerfeld[€]
\teil Eine Gitarre kostet $180{,}00\,€$. Es gibt $35\,\%$ Rabatt. Wie viel Euro Rabatt sind das? \hfill (P17) \\ \leerfeld[€]
\teil $17\,\%$ von $60\,€$ \hfill (P14) \\ \leerfeld[€]
\teil $19\,\%$ von $40\,€$ \hfill (P14) \\ \leerfeld[€]
\teil $40\,\%$ von $85\,€$ \hfill (P26F) \\ \leerfeld[€]
\teil $60\,\%$ von $45\,€$ \hfill (P26F) \\ \leerfeld[€]
\teil Bei einer Umfrage unter $800$ Haushalten hatten $35\,\%$ ein E-Bike, sechs Jahre später $52\,\%$. Um wie viele Haushalte hat die Zahl zugenommen? \hfill (P19)
\teil Eine Schule hat $650$ Schüler. Im Sommer kommen $44\,\%$ mit dem Rad, im Winter $26\,\%$. Wie viele Schüler kommen im Winter weniger mit dem Rad? \hfill (P19)
\end{teile}
\end{aufgabe}

\medskip
\uebersichtskasten{Prozentwert: erst 1 \% (Ganzes : 100), dann mal p – oder p als Dezimalzahl mal Ganzes. \\ 18 \% von 40 kg: 1 \% = 0,4 kg → 18 \% = 7,2 kg \quad 0,18 · 40 = 7,2}
